% !TEX TS-program = lualatex
% tagged-pdf-example.tex
\DocumentMetadata{
  lang        = en-US,
  pdfversion  = 2.0,
  pdfstandard = ua-2,
  % pdfstandard = a-4f, % or a-4, for archiving; see the page for details
  tagging     = on,
  tagging-setup = {math/setup=mathml-SE, math/alt/use}
}
%%%%%%%% Don't load packages you don't need, and check here for packages
%%%%%%%% not supported by the tagging project:
%%%%%%%% https://latex3.github.io/tagging-project/tagging-status/
%%%%%%%% Even amsmath and amsthm have been only partially supported!
\documentclass{article}
\usepackage{graphicx} % Required for inserting images
\usepackage[math-style=TeX]{unicode-math} % also loads fontspec
%%%%%%%% Fonts are handled differently by LuaLaTeX.
%%%%%%%% Here is an example of declaring alternative fonts.
%%%%%%%% You will need to have any fonts installed.
  % \setmathfont{XITS Math}[StylisticSet = 1]
  % \setmainfont{Crimson}

\usepackage{tikz}

\title{Tagged PDF tests}
\author{Zach Norwood}
\date{September 2025}

\usepackage{hyperref}   %%% load hyperref *last*

\begin{document}
\hypersetup{
  pdftitle={Tagged PDF tests},
  pdfauthor={Zach Norwood},
  pdfkeywords={accessibility, tagging},
  % pdfdisplaydoctitle
}
\maketitle

Give a proof by induction of the following identity, proved by other means in class.
\[
  \sum_{k=0}^n c(n,k) x^k = x^{(n)}.
\]

Matrix--vector multiplication gives a linear combination of the columns:
\[
\begin{bmatrix}
 \vert &        & \vert \\
 x_1   & \cdots & x_n   \\
 \vert &        & \vert
\end{bmatrix}
\begin{bmatrix}
 c_1 \\ \vdots \\ c_n
\end{bmatrix}
=
c_1 \begin{bmatrix} \vert \\ x_1 \\ \vert \end{bmatrix}
+ \cdots + c_n \begin{bmatrix} \vert \\ x_n \\ \vert \end{bmatrix}
\]

\section{Introduction}

\subsection{Your first theorem}

A miracle of nature:
\[
f(a) = \frac{1}{2\pi i} \oint_\gamma \frac{f(z)}{z-a}\,dz
\]

\begin{enumerate}[item-label=(\alph{*})]
%%% ^^^ no space between {enumerate} and [...]!
%%% try \Alph, \arabic, or \roman instead of \alph
\item start with A
\item then on to B
\item one more to C
\item and end with D.
\end{enumerate}

Tables require additional code to tag header rows/columns.
%%%%% full documentation here: https://texdoc.org/serve/latex-lab-table/0
\tagpdfsetup{table/header-rows={1}}
\tagpdfsetup{table/header-columns={1}}
\begin{center}
\begin{tabular}{r | l l}
 $x$ & $f(x)$ & $f(x)^2$ \\ \hline
 $1$ & $\pi$  & $\pi^2$  \\
 $2$ & $-4$   & $16$
\end{tabular}
\end{center}

Your tagging settings will persist until you change them.
So change them when you add a table that has different header properties.
\tagpdfsetup{table/header-rows={}}
\tagpdfsetup{table/header-columns={1}}
\[
\begin{array}{c| ccccc}
n   & 1 & 2 & 3 & 4 & 5 \\
\hline
p_n & 2 & 3 & 5 & 7 & 11 \\
\end{array}
\]

\clearpage

%%% Tagging of images is nontrivial.
%%% https://texdoc.org/serve/latex-lab-graphic/0

\begin{figure}[h]
\centering
\begin{tikzpicture}[style=thick, alt={The Petersen Graph}]
  \foreach \x in {0,72,2*72,3*72,4*72}
  {
  \draw (\x+18:3cm) node {} -- (\x+90:3cm) node {};
  \draw (\x+18:3cm) node {} -- (\x+18:1.7cm) node {};
  \draw (\x+18:1.7cm) node {} -- (\x+162:1.7cm) node {};
  }
\end{tikzpicture}
\caption{The Petersen Graph}
\label{petersen}
\end{figure}

%%% This needs an image file named parabola (.png, .jpg, or .pdf) in the
%%% same folder. Uncomment it once you have one, and describe your own image.
% \includegraphics[height=4cm, alt={a parabola with $y$-intercept $(0,3)$ and vertex $(3,-6)$}]{parabola}

\end{document}
